\documentclass[11pt]{article}
\usepackage[T1]{fontenc}
\usepackage[utf8]{inputenc}
\usepackage{lmodern,microtype,geometry,amsmath,amssymb,amsthm,mathtools}
\usepackage[colorlinks=true,linkcolor=blue,citecolor=blue,urlcolor=blue]{hyperref}
\geometry{margin=1in}
\hypersetup{pdftitle={Online mistake bounds for Gaussian neural networks and priors preserving their function rankings},pdfauthor={ChatGPT-6 Astra; Claude; Kaarel Hanni; Dmitry Vaintrob}}
\newtheorem{theorem}{Theorem}
\newtheorem{lemma}[theorem]{Lemma}
\newtheorem{proposition}[theorem]{Proposition}
\newtheorem{corollary}[theorem]{Corollary}
\newtheorem{remark}[theorem]{Remark}
\title{Online mistake bounds for Gaussian neural networks\\and priors preserving their function rankings}
\author{ChatGPT-6 Astra \and Claude (Anthropic) \and Kaarel H\"anni \and Dmitry Vaintrob}
\date{Research draft, 30 October 2026}
\begin{document}
\maketitle
\begin{abstract}
A small neural implementation can give polynomial learning bounds even
when its Gaussian prior probability is extremely small. This leaves a
question for online learning under distribution shift: can the learner
make only polynomially many mistakes over its entire lifetime, rather
than merely attain small test error on one distribution?

We give two results, keeping their different priors explicit. First, the
usual Gaussian tanh posterior, with global output-margin support and
exact label conditioning, has polynomial total mistakes under arbitrary
adaptive switching among finitely many fixed input environments. For
environments uniform on arbitrary subsets of the Boolean cube, the bound
is independent of learner width and polynomial in target complexity and
input dimension at fixed depth. It covers sufficiently small variance
at each width, every fixed above-mean-field exponent, and mean field.
The proof permits arbitrary probability ratios among surviving neural
functions. A small posterior tail is charged only on first appearances,
so an infinite stream incurs no horizon-times-tail penalty.

Second, we change the function prior while preserving its entire ordering
and all exact ties. Averaging weights $1/[r(r+1)]$ across tied ranks gives
a fixed prior whose target negative log probability is at most twice the
logarithm of its inclusive rank, plus a constant. Existing neural
compression and mass bounds make this logarithmic rank polynomial in
target complexity, uniformly in width. Ordinary conditioning under the
modified prior therefore gives polynomial total mistakes on every
nonanticipating adaptive input sequence. This prior is realizable on the
same neural parameter space, but is generally not an independent Gaussian.
We also give a near-minimum-energy sampling alternative that does not need
exact ties. No arbitrary-order polynomial theorem for the unchanged
large-width Gaussian learner is claimed.
\end{abstract}

\section{The distinction between ordering and probability ratios}

Consider online prediction of a fixed Boolean function: each input is
presented, a label is predicted, and the true label is then revealed.
All observations are retained. A useful simplicity principle would bound
the total number of mistakes by a polynomial in the complexity of a
target implementation, even when the queried input region changes.

Fresh uniform selection among smallest consistent circuits has such a
bound on arbitrary input orders: there are at most $s$ gate-count levels
below a target of size $s$, and each level has a logarithmic candidate
elimination budget~\cite{occam}. Real-valued neural energies do not supply
that argument. They can have many distinct levels below the target's
energy, and equal leading Gaussian rate does not imply equal function
probabilities. Exact neural symmetries do create some ties, but neither
polynomially many levels nor uniform function probabilities within every tied energy level
follows from the usual neural compression theorem.

The distinction is visible in a finite example. On points
$x_1,\ldots,x_m$, let $f$ be constantly positive and let $g_i$ be positive
except at $x_i$. Give $g_i$ energy $1+i\eta/(m+1)$ and $f$ energy
$1+\eta$, where $\eta>0$ is arbitrarily small. Exact minimum-energy
prediction makes $m$ mistakes on $x_1,\ldots,x_m$, despite the narrow
energy interval and logarithmic class size $\log(m+1)$. This is an
abstract obstruction to an inference from energy counting alone, not a
construction of these energies in the Gaussian neural model.

Our first result retains the Gaussian prior and places explicit
conditions on how inputs arrive. Our second retains the same functions,
likelihood, and prior ranking, but changes the probability ratios.
It then permits arbitrary input orders. All arguments are statistical;
we do not assert computationally efficient inference or rank evaluation.

\section{Neural model and available structural estimates}

Let $d\ge1$, $L\ge2$, and let $N$ be the common hidden width. Inputs are
$x\in\{-1,1\}^d$, augmented by a constant coordinate. There are $L-1$
hidden layers with coordinatewise tanh activation and a linear readout:
\[
 h^1(x)=\tanh(W^1(1,x)),\qquad
 h^j(x)=\tanh(W^jh^{j-1}(x)),\qquad
 z_w(x)=v^{\mathsf T}h^{L-1}(x).
\]
There are no additional internal biases. The weighted energy is
\[
 \kappa_N(w)=\frac dN\|W^1\|_F^2
       +\sum_{j=2}^{L-1}\|W^j\|_F^2+N\|v\|_2^2.
\]
For a Boolean function $g$ define its strict global-margin event and
minimum energy by
\[
 U_{N,g}=\{w:g(x)z_w(x)>1/2\text{ for every }x\},\qquad
 \kappa_N(g)=\inf_{w\in U_{N,g}}\kappa_N(w).
\]
Let $\mathcal H_N$ contain precisely the functions with nonempty
$U_{N,g}$. It is a finite set of distinct functions. Each such event is
open and has positive mass under every nondegenerate Gaussian prior.
Let $U_N=\bigcup_{g\in\mathcal H_N}U_{N,g}$ be the label-independent
global absolute-margin event.

In the strong-prior parameterization, the energy coordinates are
independent centered Gaussians of variance $\sigma^2$. Thus individual
weight variances are $\sigma^2N/d$ in the first layer, $\sigma^2$ internally,
and $\sigma^2/N$ in the readout. The finite-variance family sets
$\sigma^2=N^{-\gamma}$, giving variances
$N^{1-\gamma}/d,N^{-\gamma},N^{-\gamma-1}$. Mean field is $\gamma=1$.
Write $\mu$ for the raw Gaussian and
\[
 \pi(g)=\frac{\mu(U_{N,g})}{\mu(U_N)}
\]
for its induced function prior after global-margin conditioning. The
learning posterior further conditions on every revealed label exactly.
In particular, repeated inputs have zero posterior error.

The target complexity $C=C_L(f)$ is the least hidden width of a same-depth
tanh implementation with signed margin greater than $1/2$ everywhere
and every coefficient bounded in magnitude by that width. We assume
$C<\infty$ and allow every $N\ge C$. All polynomial claims below hold at
fixed depth; for the above-mean-field family the exponent $\gamma>1$ is
also fixed. Their degrees need not be uniform as $\gamma$ decreases to
one. The mean-field conclusion is supplied by a separate estimate.
Natural logarithms are used, and $0<\delta<1$ denotes a confidence level.

\begin{lemma}[Imported neural structural bounds and their budgets]
\label{lem:neural-inputs}
Fix depth $L\ge2$, let $C=C_L(f)<\infty$, and suppose $N\ge C$.
All function priors in this lemma use the global output margin restriction
before training. There are functions $F_L(d,J)$, $A_L(d,C)\ge1$,
$T_L(d,B)$, and $G_L(d,B)$ with polynomial majorants at fixed $L$
such that the following statements hold. The first is the norm count
\[
\log\#\{g\in\mathcal H_N:\kappa_N(g)\le J\}\le F_L(d,J)
 \quad(J\ge1),
 \qquad
 \kappa_N(f)\le B_L(d,C):=d(d+1)C^2+LC^4.
\]
For the specified finite-variance family, the raw target mass satisfies
\[
 \mu_{N,\gamma}(U_{N,f})\ge
 \exp\{-A_L(d,C)N^\gamma/2\}\qquad(\gamma\ge1).
\]
These are the compression, cloning, and target-cylinder bounds
of~\cite{strong,above,meanfield}, recorded with explicit scalar formulas
in~\cite{online}.

For strong prior, fix any $0<\beta<1$ and any admissible finite $N$.
There is $\sigma_0>0$ such that for every $0<\sigma<\sigma_0$, the class
\[
 \mathcal A=\{g\in\mathcal H_N:\kappa_N(g)\le\kappa_N(f)\}
\]
contains $f$ and satisfies
\[
 \log|\mathcal A|\le h:=F_L(d,B_L(d,C)),\qquad
 \pi(\mathcal A^c)\le\beta\pi(f).
\]
The threshold may depend on $N,d,L,f,\beta$ and the activation; no
width-uniform variance threshold is asserted. One threshold suffices for
all subsequent noiseless histories and the entire input stream.

For every fixed $\gamma>1$, put $A=A_L(d,C)$ and
\[
 B_\beta=4A+8\log(L/\beta),\qquad
 a_\beta=[8L(\sqrt{B_\beta}+1)^{L-2}]^{-1},
\]
\[
 T_\beta=\left\lceil\max\left\{
 C,[8(d+1)]^{1/(\gamma-1)},(8/a_\beta)^{2/(\gamma-1)}
 \right\}\right\rceil,
 \quad h_\beta=F_L(d,(L+3)B_\beta),
 \quad J_\beta=AT_\beta^\gamma/2.
\]
At $N\ge T_\beta$, the class
$\mathcal A=\{g:\kappa_N(g)\le(L+3)B_\beta\}$ contains $f$ and
satisfies $\log|\mathcal A|\le h_\beta$ and
$\pi(\mathcal A^c)\le\beta\pi(f)$.
At $C\le N<T_\beta$, instead $-\log\pi(f)\le J_\beta$.
Both budgets are polynomially bounded in $d,C,\log(1/\beta)$ at fixed
$L,\gamma$; their degrees may depend on $\gamma$.

For mean-field tanh, the functions $T_L,G_L$ have the following
structural meaning: for $B\ge4$ and $N\ge T_L(d,B)$, the raw total
mass of full-margin functions outside a finite class of log cardinality
at most $G_L(d,B)$ is at most $L e^{-NB/8}$.
For $\gamma=1$, use the same $A,B_\beta$ and set
\[
 T_\beta=\max\{C,T_L(d,B_\beta)\},\qquad
 h_\beta=G_L(d,B_\beta),\qquad J_\beta=AT_\beta/2.
\]
At $N\ge T_\beta$ there is a finite class $\mathcal A$ containing $f$
with $\log|\mathcal A|\le h_\beta$ and
$\pi(\mathcal A^c)\le\beta\pi(f)$; at $C\le N<T_\beta$,
$-\log\pi(f)\le J_\beta$. These budgets are polynomially bounded in
$d,C,\log(1/\beta)$ at fixed $L$.
\end{lemma}

\begin{proof}
The strong-prior assertion follows from
$2\sigma^2\log\mu_{N,\sigma}(U_{N,g})\to-\kappa_N(g)$ and
finiteness of $\mathcal H_N$: the total mass of functions with energy
strictly above the target's is $o(\pi(f))$. All energy ties stay in
$\mathcal A$; no assertion about equal weights is needed~\cite{strong,online}.
For $\gamma>1$, the spectral truncation estimate gives raw exceptional
mass at most $L\exp(-B_\beta N^\gamma/8)$ under the displayed width
conditions~\cite{above}. Dividing by the raw target mass gives at most
\[
 L\exp\{-(B_\beta/8-A/2)N^\gamma\}
 =L(\beta/L)^{N^\gamma}\le\beta.
\]
The mean-field finite-description theorem gives the identical calculation
with exceptional mass $L\exp(-NB_\beta/8)$ and width threshold
$T_L(d,B_\beta)$~\cite{meanfield}. In either case the resulting class
must contain $f$, since otherwise its relative outside mass is at least
one. At smaller widths,
$-\log\pi(f)\le-\log\mu(U_{N,f})\le AN^\gamma/2$.
The polynomial bounds and explicit quantities are also collected
in~\cite{online}. The classes and tail estimates are fixed before the
input process; exact conditioning therefore bounds posterior mass outside
each such class by $\beta$ after every realizable history.
\end{proof}

\begin{corollary}[Gaussian function ranks]
\label{cor:neural-rank}
For any prior covered by Lemma~\ref{lem:neural-inputs}, let
\[
 R_f=\#\{g\in\mathcal H_N:\pi(g)\ge\pi(f)\}.
\]
Then $\log R_f\le H_\beta$, where $H_\beta=h$ in the strong-prior
case and $H_\beta=\max\{h_\beta,J_\beta\}$ in either finite-variance
case. In particular, choosing $\beta=1/2$ gives a bound polynomial in
$d,C$ at the fixed parameters of the regime, uniformly for every
$N\ge C$. This rank is unchanged if raw global-margin masses replace
the conditioned function probabilities. The scalar $H_\beta$ bounds
both proof branches; it does not assert a common relevant class in the
smaller-width branch.
\end{corollary}

\begin{proof}
In a relevant-class branch, every $g\notin\mathcal A$ has
$\pi(g)\le\pi(\mathcal A^c)\le\beta\pi(f)<\pi(f)$.
Thus all functions counted by $R_f$, including every exact tie, belong
to $\mathcal A$, and $\log R_f\le\log|\mathcal A|$.
In a direct-mass branch, $R_f\pi(f)\le1$, so
$\log R_f\le-\log\pi(f)\le J_\beta$.
Global-margin conditioning divides all raw function masses by one common
normalizer and therefore preserves the ranking and its ties.
\end{proof}

When using strong-prior predictive limits, one takes limits of the
posterior laws after conditioning on each realizable history. This does
not require conditioning a limiting prior, which may assign the target
zero mass.

\section{Adaptive switching among fixed sources}
\label{sec:structured}

The number of energy or probability levels causes no difficulty when
fresh inputs retain sufficient randomness within each of finitely many
fixed sources. The source can change after every observation. The
following result extends the finite-class environment arguments of
\cite{occam} to a posterior with a uniformly controlled tail. Its tail
cost depends on the number of distinct possible inputs, rather than on
the duration of the stream.

Let $\mathcal X$ be a finite domain of size $m$, with fixed target $f$.
After the first $t-1$ labels, let $\Pi_{t-1}$ be a distribution on functions
that is supported on functions agreeing with all those labels and is
determined before the next input is drawn. Suppose a fixed finite set
$A$ contains $f$ and satisfies
\begin{equation}\label{eq:structured-tail}
 \log|A|\le h,\qquad \Pi_{t-1}(A^c)\le\beta
 \quad\text{after every realizable history}.
\end{equation}
For an exact consistency posterior with prior $\pi$, the sufficient
condition
\begin{equation}\label{eq:structured-prior-ratio}
 \pi(A^c)\le\beta\pi(f),\qquad \pi(f)>0,
\end{equation}
implies~\eqref{eq:structured-tail}: the posterior denominator is at least
the mass of the always-consistent target. The hypothesis is thus uniform
over rare histories as well as typical ones.

Fix distributions $D_1,\ldots,D_K$ on $\mathcal X$, with $K\ge1$.
At each round a scheduler chooses a source using the full observed past,
then obtains a fresh independent input from that source. Its choice may
also use fresh randomness, but not the unrevealed source samples or the
current private predictor draw. The learner need not observe the source
identity. Predictable mixtures of the fixed sources are included by
regarding the source choice as latent. Choose $0<\varepsilon_i\le1$ such
that
\begin{equation}\label{eq:structured-gap}
 D_i(g\ne f)\in\{0\}\cup[\varepsilon_i,1]
 \qquad(g\in A).
\end{equation}
One may use the smallest positive disagreement probability, or any lower
bound on it. If there is no positive disagreement, any such
$\varepsilon_i$ is valid.

At the actual input $X_t$, write
\[
 q_t=\Pi_{t-1}\{g:g(X_t)\ne f(X_t)\},\qquad
 S_\infty=\sum_{t\ge1}q_t.
\]
For fresh sampling from $\Pi_{t-1}$, let $M_\infty$ be the total mistakes.
The draw can occur after $X_t$ is revealed, provided its conditional law
is still $\Pi_{t-1}$.

\begin{theorem}[Posterior learning under interleaved iid sources]
\label{thm:structured-posterior}
Under~\eqref{eq:structured-tail} and~\eqref{eq:structured-gap}, arbitrary
adaptive switching among the fixed sources satisfies
\begin{equation}\label{eq:structured-expectation}
 \mathbb E S_\infty=\mathbb E M_\infty
 \le \sum_{i=1}^K
 \left[1+(h+1)\left(1+\log\frac1{\varepsilon_i}\right)\right]
       +\beta m.
\end{equation}
For $0<\delta<1$, define
\begin{gather*}
 J_i=1+\left\lceil\log_2\frac1{\varepsilon_i}\right\rceil,
 \qquad \Lambda=\sum_{i=1}^KJ_i,
 \\
 a_r=h+\log\frac{r\Lambda}{\delta},
 \qquad B_r=(\Lambda+K)a_r+2K \quad(r=2,3).
\end{gather*}
Then
\begin{align}
 \mathbb P\left\{S_\infty>
       (e-1)B_2+\log\frac2\delta+\beta m\right\}
       &\le\delta,\label{eq:structured-S-tail}\\
 \mathbb P\left\{M_\infty>
       (e-1)(B_3+\beta m)+2\log\frac3\delta\right\}
       &\le\delta.\label{eq:structured-M-tail}
\end{align}
These are separate probability statements. They hold over the entire
possibly infinite horizon and therefore also at every stopping time.
\end{theorem}

\begin{proof}
Split each realized-query error probability as
\[
 q_t=q_t^{\mathrm{in}}+q_t^{\mathrm{out}},\qquad
 q_t^{\mathrm{out}}=
 \Pi_{t-1}\{g\notin A:g(X_t)\ne f(X_t)\}.
\]
Every retained function is correct on a previously observed input.
Thus, for the actual realized queries,
\begin{equation}\label{eq:structured-tail-charge}
 q_t^{\mathrm{out}}\le
 \beta\,\mathbf1\{X_t\notin\{X_1,\ldots,X_{t-1}\}\},
 \qquad \sum_{t\ge1}q_t^{\mathrm{out}}\le\beta m.
\end{equation}
This is a pathwise inequality and needs no assumption on input order.
It is different from a bound on risk averaged over the next source draw.

For the inside contribution, couple source $i$ to its own independent iid
stream. Let $R_{i,n}$ be the maximum $D_i$-risk of a member of $A$ agreeing
with the first $n$ labels in that stream, ignoring labels from other
sources. A function with risk $u$ survives with probability $(1-u)^n$;
therefore, for $n\ge1$,
\begin{equation}\label{eq:structured-prefix-tail}
 \mathbb P\{R_{i,n}>u\}\le\min\{1,e^{h-nu}\}.
\end{equation}
At a visit to source $i$ after $n$ of its samples have been seen, all
inside-$A$ posterior functions survive that local prefix. Consequently,
conditional on the observed past and the chosen source, but \emph{before}
the new input is drawn,
\begin{equation}\label{eq:structured-inside-drift}
 \mathbb E[q_t^{\mathrm{in}}\mid\text{past, source }i]
 \le R_{i,n}.
\end{equation}
The same bound holds for the indicator of a sampled mistake from a
function in $A$. Arbitrary relative posterior weights within $A$ are
allowed. Other sources' observations can only further reduce the set of
possible functions.

The random variable $R_{i,n}$ is either zero or at least
$\varepsilon_i$. For $u>0$,
\[
 \sum_{n\ge1}\min\{1,e^{h-nu}\}
 \le\int_0^\infty\min\{1,e^{h-tu}\}\,dt=\frac{h+1}{u}.
\]
For $0<u<\varepsilon_i$, use instead
$\mathbb P(R_{i,n}>u)\le\min\{1,e^{h-n\varepsilon_i}\}$,
by union bounding all positive-risk functions. Integrating the tail and
using Tonelli's theorem gives
\[
 \sum_{n\ge1}\mathbb E R_{i,n}
 \le(h+1)\left(1+\log\frac1{\varepsilon_i}\right).
\]
The initial local round costs at most one. Each local counter $(i,n)$ is
used at most once, irrespective of the adaptive interleaving. Summing
\eqref{eq:structured-inside-drift}, adding
\eqref{eq:structured-tail-charge}, and using conditional expectation for
the fresh sampled predictions proves~\eqref{eq:structured-expectation}.

We give the high-probability argument with a generic $a>0$. For
$1\le j\le J_i$, set $n_{i,j}=\lceil a2^j\rceil$. Except on an event of
probability at most $\Lambda e^{h-a}$, all prefix inequalities
\[
 R_{i,n_{i,j}}\le2^{-j}
\]
hold. Since $2^{-J_i}<\varepsilon_i$, the last one forces risk zero.
Monotonicity of $R_{i,n}$ supplies a deterministic envelope: use one
before $n_{i,1}$, use $2^{-j}$ from $n_{i,j}$ until $n_{i,j+1}-1$, and
use zero from $n_{i,J_i}$ onwards. Its sum over all local counts for
source $i$ is at most
\[
 2a+1+\sum_{j=1}^{J_i-1}(a+2^{-j})
 \le(J_i+1)a+2.
\]
Stop the inside process before using any prefix that violates this
applicable envelope. Whether the current prefix violates it is known
from the observed local history before drawing the next input. Until
stopping, the sum of conditional mean bounds is at most
$B=(\Lambda+K)a+2K$ on every path. For either the stopped fractional
variable $q_t^{\mathrm{in}}$ or the stopped inside-mistake indicator,
$0\le Y_t\le1$ and
\[
 \mathbb E[e^{Y_t}\mid\text{past, chosen source}]
 \le 1+(e-1)\mathbb E[Y_t\mid\text{past, chosen source}].
\]
The exponential-supermartingale argument therefore bounds its total by
$(e-1)B+\log(1/\eta)$ apart from probability $\eta$. This holds uniformly
in time. On the simultaneous good-prefix event, stopping never occurs.
Taking $a=a_2$, $\eta=\delta/2$, and adding the pathwise outside sum proves
\eqref{eq:structured-S-tail}.

For sampled mistakes use $a=a_3$ and $\eta=\delta/3$ for the inside
process. For the outside process a different conditioning stage is
appropriate: condition on the past, the chosen source, and the
\emph{realized current input}, immediately before the private function
draw. The outside-mistake indicator $Z_t^{\mathrm{out}}$ then has
conditional mean $q_t^{\mathrm{out}}$, and
\[
 \mathbb E\!\left[
 e^{Z_t^{\mathrm{out}}-(e-1)q_t^{\mathrm{out}}}
 \mid\text{past, source, current input}\right]\le1.
\]
Taking the further expectation over the source and input makes the
product a nonnegative supermartingale at the ends of complete rounds.
By~\eqref{eq:structured-tail-charge}, its compensator is bounded
pathwise by $(e-1)\beta m$. Thus outside mistakes total at most
$(e-1)\beta m+\log(3/\delta)$, apart from probability $\delta/3$.
A union bound over prefix failure, the inside tail, and the outside tail
proves~\eqref{eq:structured-M-tail}. Infinite sums follow by monotone
limits, or directly from the time-uniform supermartingale bounds.
\end{proof}

Pointwise, the number of posterior-majority mistakes is at most
$2S_\infty$, and the number of failures to assign the correct label
posterior probability strictly greater than $2/3$ is at most
$3S_\infty$. Thus~\eqref{eq:structured-expectation} gives these expectation
bounds with factors two and three, respectively;
\eqref{eq:structured-S-tail} gives their high-probability counterparts.
Majority prediction may depend on the query; it is controlled through
this posterior-risk comparison, not by assuming the majority function
itself belongs to $A$.

\paragraph{Finite pools.}
Fix $K$ finite pools of sizes $m_i\ge1$, independently uniformly permute
each, and adaptively interleave their unrevealed entries. Repeated inputs
within or between pools are allowed. The random-pool prefix argument of
\cite{occam}, applied to the inside posterior, gives
\begin{equation}\label{eq:structured-pools}
 \mathbb E M\le
 K+2(h+1)\sum_{i=1}^K H_{m_i-1}+\beta m,
 \qquad H_0=0.
\end{equation}
Indeed, after $n$ entries of pool $i$, the conditional inside risk is at
most $m_iR_{i,n}/(m_i-n)$, and
$\sum_{n=1}^{m_i-1}m_i/[n(m_i-n)]=2H_{m_i-1}$. The outside term is still
\eqref{eq:structured-tail-charge}, regardless of duplicate entries.
For pools with at most $2^d$ entries, this is
$O(K(h+1)(d+1))+\beta2^d$.
The same stopping and tail split gives a sampled high-probability bound
$(e-1)(B+\beta m)+2\log(3/\delta)$, using
$a=h+\log(3\sum_i m_i/\delta)$ and
$B=K+2a\sum_i H_{m_i-1}$. For cumulative posterior risk use the
corresponding two-event allocation, as in~\eqref{eq:structured-S-tail}.

\begin{corollary}[Structured shifts for the original Gaussian NNBayes learners]
\label{cor:structured-neural}
Use any of the strong-prior, fixed-$\gamma>1$, or mean-field tanh regimes
of Lemma~\ref{lem:neural-inputs}, with the same target and depth
qualifications. Let $m=2^d$, $\beta=1/m$, and let $H_\beta$ majorize both
the lemma's relevant-class logarithmic count $h_\beta$ and its direct
prior-cost budget $J_\beta$ where that alternative is used. For $K\ge1$
fixed sources uniform on arbitrary nonempty subsets of the cube,
\begin{equation}\label{eq:structured-neural-expectation}
 \mathbb E M_\infty\le
 K\bigl[1+(H_\beta+1)(1+d\log2)\bigr]+1.
\end{equation}
With
\[
 a=H_\beta+\log\frac{3K(d+1)}\delta,
 \qquad B=K(d+2)a+2K,
\]
one also has
\begin{equation}\label{eq:structured-neural-tail}
 \mathbb P\left\{M_\infty>
 (e-1)(B+1)+2\log\frac3\delta\right\}\le\delta.
\end{equation}
Thus at fixed depth, and fixed $\gamma$ where applicable, these are
polynomial total-mistake bounds in $d,C,K,\log(1/\delta)$, independently
of learner width. The scheduler may switch indefinitely and arbitrarily
abruptly among the fixed sources. Corresponding majority and confidence
bounds follow from the cumulative posterior risk. Fixed finite pools
with at most $2^d$ entries each have the same polynomial conclusion.
\end{corollary}

\begin{proof}
In the relevant-class branch of Lemma~\ref{lem:neural-inputs}, the
prior-ratio condition gives~\eqref{eq:structured-tail} after every
history. Every positive disagreement probability on a uniform subset of
the cube is at least $2^{-d}$, so use $\varepsilon_i=2^{-d}$ for every
source. Then $J_i=d+1$, $\Lambda=K(d+1)$, $h\le H_\beta$, and
$\beta m=1$. Theorem~\ref{thm:structured-posterior} gives the displayed
bounds.

In the direct-prior-cost branch, no common finite class is needed or
asserted. Exact Bayesian telescoping gives
$S_\infty\le-\log\pi(f)\le J_\beta\le H_\beta$ for every input sequence,
and posterior sampling has tail
$(e-1)H_\beta+\log(1/\delta)$. These are bounded by the displayed budgets.
The same branch distinction gives the posterior-risk claim. For finite
pools take the maximum of the direct budget and the relevant-class
budget in~\eqref{eq:structured-pools}; no finite-class assertion is made
in the direct branch. The imported budgets are polynomial after substituting
$\log(1/\beta)=d\log2$.
\end{proof}

For strong prior the variance is chosen sufficiently small at each fixed
width before the stream begins, with the original possibility that its
threshold depends on width. One such choice covers all histories and
all subsequent switches. Any simultaneous limit of the correctly
conditioned predictive posteriors as the variance tends to zero at that
fixed width also satisfies the result: each limiting posterior is
supported on the target-energy sublevel, giving $\beta=0$. This refers
to limits \emph{after conditioning on each history}; conditioning a
limiting prior that has assigned zero mass to the target is not the
argument.

The sources themselves must be fixed in advance. A polynomial number of
sources with arbitrary switching is a substantive extension of the
uniform-stream theorem, but does not cover every deterministic order:
encoding an arbitrary cube order as singleton sources uses $K=2^d$.

\section{Rank flattening of the function prior}
\label{rank:section}

All logarithms in this section are natural. Let $\mathcal H$ be a finite
set of distinct Boolean functions, let $M=|\mathcal H|\ge1$, and let
$\pi$ be a strictly positive prior on $\mathcal H$. For $f\in\mathcal H$,
define its strict and inclusive ranks by
\[
 a_f=|\{g:\pi(g)>\pi(f)\}|,\qquad
 b_f=|\{g:\pi(g)\ge\pi(f)\}|.
\]
Define the \emph{rank-flattened prior}
\begin{equation}\label{rank:prior}
 Z_M=\frac{M}{M+1},\qquad
 q_\pi(f)=\frac{1}{Z_M(a_f+1)(b_f+1)}.
\end{equation}
This transformation is performed once, before observing any labels.
Afterwards, $q_\pi$ is updated by ordinary conditioning on consistency.

\begin{proposition}[Normalization, preferences, and target cost]
\label{rank:normalization}
The law $q_\pi$ is a probability distribution and preserves every strict
ordering and every tie of $\pi$:
\[
 \pi(f)>\pi(g)\iff q_\pi(f)>q_\pi(g),\qquad
 \pi(f)=\pi(g)\iff q_\pi(f)=q_\pi(g).
\]
In particular, the two consistency posteriors have the same sets of
maximum-a-posteriori functions on every history. Writing $R_f=b_f$,
\begin{equation}\label{rank:cost}
 \log\frac1{q_\pi(f)}
 =\log Z_M+\log(a_f+1)+\log(b_f+1)
 \le\log\bigl(R_f(R_f+1)\bigr)
 \le2\log(R_f+1).
\end{equation}
\end{proposition}

\begin{proof}
Order functions by decreasing $\pi$-mass. A tie block occupies consecutive
ranks $a+1,\ldots,b$. Its mean weight under $w_r=1/[r(r+1)]$ is
\[
 \frac1{b-a}\sum_{r=a+1}^b\frac1{r(r+1)}
 =\frac1{b-a}\left(\frac1{a+1}-\frac1{b+1}\right)
 =\frac1{(a+1)(b+1)}.
\]
Thus the unnormalized masses in~\eqref{rank:prior} sum to
$\sum_{r=1}^M w_r=Z_M$. Since $w_r$ is strictly decreasing, its averages
over successive disjoint rank blocks strictly decrease. This proves the
ordering and tie assertions. Conditioning multiplies all surviving
masses by a common positive constant, so it preserves the maximizing
set. Finally $a_f+1\le b_f=R_f$ and $Z_M<1$ give~\eqref{rank:cost}.
\end{proof}

\begin{theorem}[Arbitrary-order prediction after rank flattening]
\label{rank:online}
Fix a target $f_*\in\mathcal H$. At round $t$, an input $x_t$ is chosen,
the learner predicts, and the label $f_*(x_t)$ is revealed. The inputs may
depend on the whole preceding history, including past randomized
predictions, but are chosen before the learner's fresh randomization at
the current round. Let $q_t$ be the consistency posterior of the fixed
prior $q_\pi$, and set
\[
 p_t=q_t\{g:g(x_t)=f_*(x_t)\},\qquad
 B_* =\log\frac1{q_\pi(f_*)}.
\]
Uniformly over such input-selection rules, including infinite streams,
the following hold.
\begin{enumerate}
\item Fresh posterior sampling has expected total mistakes at most
$B_*$, which is bounded by~\eqref{rank:cost}.
\item Posterior majority, with any tie rule, has at most $B_*/\log2$
mistakes.
\item The number of rounds with $p_t\le2/3$ is at most
$B_*/\log(3/2)$.
\end{enumerate}
For fresh posterior sampling, with probability at least $1-\delta$, the
total number of mistakes is also at most
$(e-1)B_*+\log(1/\delta)$, for every $0<\delta<1$.
\end{theorem}

\begin{proof}
Let $V_t$ be the functions consistent with the first $t-1$ labels.
Because $f_*\in V_t$, all posterior denominators are positive. On every
realized path and at every finite horizon $T$,
\begin{equation}\label{rank:potential}
 \prod_{t=1}^T p_t=q_\pi(V_{T+1})\ge q_\pi(f_*),
 \qquad \sum_{t=1}^T-\log p_t\le B_*.
\end{equation}
The conditional mistake probability of a fresh posterior draw is
$u_t=1-p_t\le-\log p_t$. Taking expectations and then applying monotone
convergence proves the first assertion, including the infinite horizon.
A majority mistake implies $p_t\le1/2$, and a round with $p_t\le2/3$
contributes at least $\log(3/2)$ to~\eqref{rank:potential}. This proves
the next two assertions pathwise.

For the last assertion, let $I_t$ be the sampled mistake indicator.
Conditionally on the history and $x_t$,
\[
 \mathbb E[e^{I_t}]=1+(e-1)u_t\le e^{(e-1)u_t}.
\]
Consequently $\exp(\sum_{t\le T}I_t-(e-1)\sum_{t\le T}u_t)$ is a
nonnegative supermartingale starting at one. Since
$\sum_{t\le T}u_t\le B_*$ pathwise, Markov's inequality gives the stated
bound at each finite horizon. Letting the horizon increase proves the
same bound for the total number of mistakes.
\end{proof}

The theorem concerns posterior sampling and posterior majority, not
posterior MAP. In fact Proposition~\ref{rank:normalization} shows that
rank flattening leaves the MAP prediction rule unchanged if its tie
rule is unchanged. Nor does the theorem concern repeatedly rank
flattening a changing posterior: its proof uses one fixed predata prior.

\begin{corollary}[Arbitrary-order neural learning with preserved function rankings]
\label{rank:neural}
Use any Gaussian regime and eligible width in
Lemma~\ref{lem:neural-inputs}, and transform its actual induced function
prior by~\eqref{rank:prior}. With the rank budget $H_{1/2}$ from
Corollary~\ref{cor:neural-rank}, the online guarantees of
Theorem~\ref{rank:online} hold with
\[
 B_*\le 2H_{1/2}+\log2.
\]
This is polynomial in $d,C$ at fixed depth, and fixed $\gamma>1$ where
applicable, independently of learner width. It includes the positive
mean-field variance. In the strong-prior case it holds for every
sufficiently small positive variance at each fixed width, with the
variance threshold independent of the stream and its length.
\end{corollary}
\begin{proof}
Corollary~\ref{cor:neural-rank} gives $\log R_f\le H_{1/2}$.
Use~\eqref{rank:cost} and
$\log[R(R+1)]\le2\log R+\log2$, then apply
Theorem~\ref{rank:online}. Neither the construction of the transformed
prior nor its posterior updates require knowing the target or its
complexity. That complexity is used only to prove the bound.
\end{proof}

An alternative ranks directly by increasing minimum energy
$\kappa_N(g)$, treating equal energies as ties. Its inclusive target rank
is at most the energy-sublevel count from Lemma~\ref{lem:neural-inputs}.
The same construction then gives an arbitrary-order budget
$2F_L(d,B_L(d,C))+\log2$ at every $N\ge C$, without a variance parameter.
It preserves energy ordering, rather than exact finite-variance Gaussian
probability ordering. No limiting equality of prefactors is assumed.

\paragraph{Realization as a parameter prior.}
Let $\mu^U$ be the original parameter law conditioned on the global
margin event, and let $g(w)\in\mathcal H$ be its induced function, so
$\mu^U\{w:g(w)=f\}=\pi(f)$. The probability law
\begin{equation}\label{rank:lift}
 \frac{d\nu}{d\mu^U}(w)=\frac{q_\pi(g(w))}{\pi(g(w))}
\end{equation}
has induced function law $q_\pi$. Indeed, its integral over the event
$\{g(w)=f\}$ is $q_\pi(f)$. It preserves the conditional distribution
of parameters within each function event. This is a global reweighting
of function masses; it generally destroys Gaussianity and coordinate
independence. It is neither ordinary Gaussian temperature adjustment nor
an efficient sampling claim. Raising a function mass to a power and
raising a parameter density to that power are different operations,
since integration and taking powers do not generally commute.

\paragraph{An alternative that samples near the minimum energy.}
Exact real-valued energies need not be tied for a related construction
to work. Let $E:\mathcal H\to[0,\infty)$ and
$H_j=\{g:E(g)\le2^j\}$ for $j\ge0$. On each history choose the least $j$
for which $H_j$ has a consistent function, and sample uniformly over its
distinct consistent functions. Retain every past label when the budget
increases. If the least surviving energy is $m\ge1$, every sampled energy
is at most $2m$; for $m<1$, it is at most one.

\begin{corollary}[Geometric energy bands]\label{rank:bands}
For the preceding learner, put
$J_* =\lceil\log_2\max\{1,E(f_*)\}\rceil$. Then
\[
 \mathbb E[\text{total mistakes}]
 \le J_*+\sum_{j=0}^{J_*}\log\max\{1,|H_j|\}.
\]
In particular, if $\log|H_j|\le A2^{jp}$ whenever $H_j$ is nonempty,
with $A>0$ and $p>0$, this bound is at most
\[
 J_*+A\frac{2^{p(J_*+1)}-1}{2^p-1},
\]
and hence polynomial in the target energy for fixed $p$. The neural energy
and entropy bounds of Lemma~\ref{lem:neural-inputs} give a
width-independent polynomial at fixed depth.
\end{corollary}

\begin{proof}
The selected budget never decreases and never exceeds $J_*$. During a
fixed-budget phase, let $n_t$ be its number of consistent functions.
Except on a round that empties the phase, the error probability obeys
$1-n_{t+1}/n_t\le\log(n_t/n_{t+1})$, so these costs telescope to at most
$\log|H_j|$. Each emptied phase costs at most one additional mistake;
there are at most $J_*$ such phases because $H_{J_*}$ contains the target.
Sum the bounds and, for an infinite stream, use monotone convergence.
\end{proof}

A fixed Bayesian alternative mixes the uniform laws on nonempty $H_j$
with weights proportional to $2^{-j-1}$. Its target negative log-mass is
at most $\log|H_{J_*}|+(J_*+1)\log2$, so the potential
argument~\eqref{rank:potential} immediately gives the corresponding online
bounds. Both energy constructions introduce ties between different
energies in one geometric bin; rank flattening preserves every original
strict probability preference.


\section{What changes and what remains unresolved}

The structured-environment theorem concerns the original Gaussian
posterior. It never replaces unequal posterior weights by uniform ones.
The sources are fixed before the process, while switches or mixture
weights can depend on the full observed past. Allowing an arbitrary new
source at no cost would invalidate the conclusion: encoding every
adversarial query as its own singleton source makes their number as large
as the domain. The theorem therefore supplies a concrete class of
distribution shifts, rather than a claim that most imaginable sequences
are benign.

The prior transformation concerns a different learner. It preserves which
function is preferred to which, but changes how strongly it is preferred.
This is enough to improve the general online bound because the
logarithm of the number of functions preferred to a small target is
polynomial, even when the original target probability is extremely small.
Since the maximizing functions do not change, the improvement is a
theorem about posterior sampling and posterior majority, not about
choosing one posterior mode.

Function-level flattening also differs from increasing the variance of
the original independent Gaussian weights. Function masses integrate
parameter density over different regions of parameter space; changing a
scalar temperature does not generally remove their volume differences.
No theorem here identifies either function-level construction with an
independent-coordinate weight prior. Nor does counting low-energy
functions alone establish their Gaussian mass ratios.

For the unchanged Gaussian learner at unrestricted large width, a
polynomial total-mistake bound for every adaptive input order remains
open in the fixed-depth setting considered here. Its elementary
bound $-\log\pi(f)$ can grow with width or diverge as variance decreases;
this fact alone is not a lower bound on actual mistakes. The geometric
bands and rank-based prior are precise positive alternatives, not a
resolution of that remaining question. The global-margin support
condition and fixed-depth quantifiers remain part of all neural
applications in this paper.

\paragraph{Verification.}
The new online transfers and function-prior constructions have
conventional proofs and independent AI-agent mathematical reviews.
They are not Lean formalizations. Existing formal verification of the
imported neural learning results does not automatically certify these
new consequences.

\begin{thebibliography}{9}
\raggedright
\bibitem{strong} ChatGPT-6 Astra, Claude Fable 5 and Opus 5, Kaarel H\"anni, and Dmitry Vaintrob.
\emph{Width-independent generalization in the strong-prior limit of Bayesian neural networks}.
\url{https://github.com/kaarelh/AI-works/tree/main/neural-network-generalization/main-results/strong-prior-generalization}.
\bibitem{above} ChatGPT-6 Astra, Claude, Kaarel H\"anni, and Dmitry Vaintrob.
\emph{Width-uniform Bayesian generalization above the mean-field scaling}.
\url{https://github.com/kaarelh/AI-works/tree/main/neural-network-generalization/main-results/gamma-above-one}.
\bibitem{meanfield} ChatGPT-6 Astra, Kaarel H\"anni, and Dmitry Vaintrob.
\emph{Mean-field Bayesian circuit learning at every width above target complexity}.
\url{https://github.com/kaarelh/AI-works/tree/main/neural-network-generalization/main-results/meanfield-all-width}.
\bibitem{online} ChatGPT-6 Astra, Claude, Kaarel H\"anni, and Dmitry Vaintrob.
\emph{Polynomially many total mistakes for Bayesian neural-network learning}.
\url{https://github.com/kaarelh/AI-works/tree/main/neural-network-generalization/works-in-progress/online-mistakes/upper}.
\bibitem{occam} ChatGPT-6 Astra, Claude, Kaarel H\"anni, and Dmitry Vaintrob.
\emph{Exponential mistakes and polynomial guarantees for circuit Occam learning}.
\url{https://github.com/kaarelh/AI-works/tree/main/neural-network-generalization/works-in-progress/online-mistakes/occam-distribution-shift}.
\end{thebibliography}
\end{document}
